\documentclass[11pt,a4paper]{article}
\usepackage[T1]{fontenc}
\usepackage{lmodern,amsmath,amssymb,amsthm,mathtools,microtype}
\usepackage[margin=27mm]{geometry}
\usepackage[colorlinks=true,linkcolor=blue,citecolor=blue,urlcolor=blue]{hyperref}
\newtheorem{theorem}{Theorem}[section]
\newtheorem{lemma}[theorem]{Lemma}
\newtheorem{proposition}[theorem]{Proposition}
\newtheorem{corollary}[theorem]{Corollary}
\newcommand{\Q}{\mathbb Q}
\newcommand{\Z}{\mathbb Z}
\newcommand{\hF}{h_{\mathrm F}}
\DeclareMathOperator{\Tr}{Tr}
\DeclareMathOperator{\ord}{ord}
\setlength{\emergencystretch}{2em}
\title{Multiplicity of modular factors and quadratic contact-height bounds}
\author{}\date{10 November 2026}
\begin{document}\maketitle
\begin{abstract}
Under the normalized modular realization of the five-contact construction,
we prove a contact-height estimate with exponent \(2+\eta\) for every
\(\eta>0\). The estimate is uniform over all rational arguments outside
the contacts. The coefficient field of a modular form contains its
nebentypus values, and all coefficient conjugates of the form have the
same restriction over the field realizing the quaternionic surface.
Consequently the corresponding modular factor is isogenous to a power
of that surface. Superadditivity of stable Faltings height along a
filtration supplies one surface of height at most the ambient height
divided by the multiplicity. A dimension-two isogeny estimate then
transfers this bound to the original surface. This removes the
character-order factor without imposing congruence restrictions on
the contact primes.
\end{abstract}

\section{A height estimate using multiplicity}

All Faltings heights use the original stable normalization. We use
R\'emond's inequality
\begin{equation}\label{eq:remond}
 \hF(B)+\hF(J/B)\le\hF(J)
\end{equation}
for an abelian subvariety \(B\subset J\), and Bost's lower bound
\(\hF(C)\ge-3\dim(C)/2\). The first is
\cite[Corollary~1.2 and the paragraph after Corollary~1.3]{Remond};
Bost's bound in this normalization is recalled in the same reference,
with the sharper constant $\log(\pi\sqrt2)<3/2$. Both assertions hold
after arbitrary finite extension.

\begin{lemma}\label{lem:average}
Let \(L\) be a number field, let \(J/L\) have dimension \(g\), and
let \(A/L\) be an abelian surface. Suppose \(J\) contains an
\(L\)-abelian subvariety isogenous over \(L\) to \(A^r\), with \(r\ge1\).
Put \(d_L=[L:\Q]\) and
\[
 H_J=\max(1,\hF(J)+3g/2).
\]
Then
\[
 \hF(A)\le \frac{H_J}{r}
 +8\left(32\log14+\log d_L+
       \log\max\left(1,\log d_L,\frac{H_J}{r}\right)\right).
\]
\end{lemma}
\begin{proof}
Choose an \(L\)-isogeny \(A^r\to B\subset J\), and let \(B_i\)
be the image of its first \(i\) factors. Each quotient
\(C_i=B_i/B_{i-1}\) is an abelian surface \(L\)-isogenous to \(A\).
Applying \eqref{eq:remond} repeatedly gives
\[
 \sum_{i=1}^r\hF(C_i)\le\hF(B)
 \le\hF(J)+\tfrac32(g-2r).
\]
Thus at least one \(C_i\) has height at most \(H_J/r\).
The dimension-two specialization of
\cite[Theorem~1.9(1)]{GaudronRemond} supplies an \(L\)-isogeny
from this \(C_i\) to \(A\) of degree at most
\[
 \left(14^{32}d_L
       \max(1,\log d_L,\hF(C_i))\right)^{16}.
\]
Stable Faltings height changes by at most half the logarithm of
the degree in the required direction. Substituting the upper bound
for \(\hF(C_i)\) proves the result.
\end{proof}

The lemma does not identify heights of isogenous varieties.
It first selects an actual surface subquotient of small height and
then applies the isogeny estimate only in dimension two. In particular,
no isogeny estimate in dimension \(2r\) is required.

\section{Multiplicity from rational restricted traces}

Let \(f\) be a primitive weight-two newform of level \(N\), nebentypus
\(\varepsilon\), and coefficient field
\(E_f=\Q(a_n(f):n\ge1)\). Write \(e=[E_f:\Q]\).
Let \(A_f/\Q\) be its modular abelian variety, of dimension \(e\).
For a prime \(\ell\), its Tate module over \(\overline{\Q}_\ell\)
is the direct sum of the two-dimensional representations \(r_f^\tau\)
over the \(e\) embeddings of \(E_f\). These facts are the modular
construction recalled in \cite[Section~3]{Ribet}.

\begin{proposition}\label{prop:multiplicity}
Let \(L/\Q\) be a number field and let \(A/L\) be an absolutely
simple abelian surface. Suppose
\[
 V_\ell(A)\otimes\overline{\Q}_\ell
       \simeq (r_f|_{G_L})^{\oplus2},
\]
where \(r_f|_{G_L}\) is absolutely irreducible, and suppose this
identification belongs to the compatible systems at all good places.
Then \(e\) is even and
\[
                         (A_f)_L\sim_L A^{e/2}.
\]
If \(s=\ord(\varepsilon)\), its multiplicity \(r=e/2\) satisfies
\[
                     r\ge \varphi(s)/2.
\]
If \(s\) is a power of two, including \(s=1\), then \(s/r\le4\).
\end{proposition}
\begin{proof}
At a good place \(v\) of \(L\), the compatible-system identification gives
\[
 \Tr r_f(\operatorname{Frob}_v)
       =\tfrac12\Tr V_\ell(A)(\operatorname{Frob}_v)\in\Q.
\]
The left side is an algebraic number in \(E_f\), independently of
the chosen coefficient embedding. Indeed, outside the finitely
many bad rational primes, it is the trace of the appropriate
power of a rational-prime Frobenius, hence a polynomial in
\(a_p(f)\) and \(p\varepsilon(p)\).
Every embedding \(\tau:E_f\hookrightarrow\overline{\Q}_\ell\)
therefore gives the same restricted Frobenius traces.
Chebotarev density and semisimplicity imply
\[
                         r_f^\tau|_{G_L}\simeq r_f|_{G_L}.
\]
Consequently
\[
 V_\ell((A_f)_L^2)\otimes\overline{\Q}_\ell
       \simeq V_\ell(A^e)\otimes\overline{\Q}_\ell.
\]
The semisimple characters already agree over \(\Q_\ell\).
Faltings' isogeny theorem~\cite{Faltings} gives \((A_f)_L^2\sim_L A^e\).
As \(A\) is absolutely simple, every simple isogeny factor of
\((A_f)_L\) is \(L\)-isogenous to \(A\).
Thus \((A_f)_L\sim_L A^r\), and dimension comparison gives \(2r=e\).
This argument also proves the parity assertion; it does not assume it.

For \(p\nmid N\), the Hecke relation
\[
                  a_p(f)^2-a_{p^2}(f)=p\varepsilon(p)
\]
shows that \(\varepsilon(p)\in E_f\). Primes in residue classes
prime to \(N\) supply the full image of \(\varepsilon\), so
\(\Q(\zeta_s)\subset E_f\). Hence \(e\ge\varphi(s)\).
For \(s\ge4\) a power of two, \(\varphi(s)=s/2\), giving
\(r\ge s/4\). For \(s=1,2\), the already proved \(r\ge1\)
gives the same stated bound.
\end{proof}

Put \(H=\ker\varepsilon\). Since the weight is two,
$\varepsilon(-1)=1$, so $-1\in H$, as required for the modular-degree
formula below. The modular variety \(A_f\) is a
quotient of \(J_H=\operatorname{Jac}(X_1(N)/H)\): every conjugate
of \(f\) has the same kernel of its nebentypus, so all its
differentials descend to this diamond quotient. Poincar\'e
reducibility therefore supplies a subvariety of \(J_H\) isogenous
to \(A_f\). Proposition~\ref{prop:multiplicity} and
Lemma~\ref{lem:average} apply to that subvariety over \(L\).

\section{The normalized arithmetic input}

Let
\[
 B_*=\prod_{p\le76176}p,\qquad C_0=210^{300},\qquad
 \mathcal D=\{0,1,2401,-343/32\}.
\]
For \(a=u/v\in\Q\setminus\mathcal D\), in lowest terms with \(v>0\),
write \(h(a)=\log\max(|u|,v)\). Let \(M=M_5(a)\) be the product
of primes greater than \(76176\) occurring to a positive odd
power in at least one of
\[
                 v,\ u,\ u-v,\ u-2401v,\ 32u+343v.
\]
Thus \(M\) is squarefree and coprime to \(B_*\). Put \(R=B_*M\)
and \(\mathcal E(R)=\prod_{p\mid R}(1+1/p)\).

We retain the full normalized modular realization in
\cite[proof of \texttt{o8-height}, equations (LR2)--(LR5)]{Source}.
Precisely, there is a surface \(A\) with
\[
                  h(a)\le2\max(0,\hF(A))+128.
\]
In the CM case its height is less than \(854R\).
Otherwise \(A\) is absolutely simple with geometric endomorphism
algebra the quaternion division algebra of discriminant \(210\).
The normalized representation has the form
\[
 \rho_\ell=b^{-1}a_\sigma\mu_\sigma\circ\sigma
              =r_f\oplus r_f,
\]
and restricts to \(V_\ell(A)\) over the normal kernel field \(L'\)
of \(b\) on the original fiber field. All geometric endomorphisms of $A$ are defined over $L'$. Faltings'
endomorphism theorem identifies the commutant of its Tate representation
over $\overline{\Q}_\ell$ with $M_2(\overline{\Q}_\ell)$. Thus the
repeated two-dimensional constituent is absolutely irreducible. The
modular representations are compatible with the Tate module over $L'$,
as required in Proposition~\ref{prop:multiplicity}. The source proves
\[
 N\le C_0R^2,\qquad \operatorname{rad}(N)\mid R.
\]
The minimal local characters and determinant calculation in
\cite[Sections~1--2]{Characters} give
\[
 [L':\Q]\le1152U(R),\qquad s=\ord\varepsilon\mid U(R),
\]
where
\[
 U(R)=\max\bigl(\{2\}\cup
             \{2^{v_2(p-1)}:p\mid R,\ p\text{ odd}\}\bigr).
\]
In particular \(s\) is a power of two,
\(U(R)\mid\varphi(R)\), and \(U(R)\mathcal E(R)<R\).
No bound on \(U(R)\) independent of \(R\) is assumed.

\section{The quadratic contact-height estimate}

\begin{theorem}\label{thm:height}
Under the normalized arithmetic input above, one has
\[
 h(a)\le67C_0(B_*M)^2\mathcal E(B_*M)\log(B_*M).
\]
Consequently, with
\[
               C_2=67C_0B_*^2\mathcal E(B_*)\log B_*,
\]
one has
\begin{equation}\label{eq:quadratic}
                    h(a)\le C_2 M^2(1+\log M)^2.
\end{equation}
For every \(\eta>0\) there is an effective constant \(C_\eta\)
such that \(h(a)\le C_\eta M^{2+\eta}\).
The exact integer certificate gives $C_2<2^{221231}$. In particular
\[
                         h(a)<7C_2 M^{7/3}<2^{221234}M^{7/3}.
\]
\end{theorem}
\begin{proof}
Put \(\ell_R=\log R\) and \(\mathcal T=C_0R^2\mathcal E(R)\).
The modular degree is
\[
 d_H=N s\prod_{p\mid N}(1+1/p)
       \le C_0R^2s\mathcal E(R)<C_0R^3.
\]
The unit-fiber differential-evaluation estimate
\cite[Section~3]{Characters} and its explicit numerical comparison give
\[
 H_J=\max(1,\hF(J_H)+3\dim J_H/2)\le8d_H\ell_R.
\]
Here the exact inequalities \(C_0^2<B_*\), \(144<B_*\), and
\(\ell_R\ge1\) give
\(\log N\le5\ell_R/2\), \(\log d_H\le7\ell_R/2\), and
\(\log12\le\ell_R/2\); the genus is at most \(d_H\).

By Proposition~\ref{prop:multiplicity}, \(J_H\) contains over \(L'\)
a subvariety isogenous to \(A^r\), with \(s/r\le4\). Thus
\[
                            H_J/r\le32\mathcal T\ell_R.
\]
Write \(d'=[L':\Q]\). Since \(d'\le1152R\),
\(\log d'<9\ell_R\). Also
\[
 \log(32\mathcal T\ell_R)<9\ell_R,\qquad
 32\mathcal T\ell_R>\max(1,\log d').
\]
Indeed \(\mathcal E(R)\le R\), \(\log C_0<\ell_R/2\),
\(\log\ell_R\le\ell_R\), and \(\log32<4\le4\ell_R\).
Lemma~\ref{lem:average} and \(\log14<3\) give
\[
 \hF(A)\le32\mathcal T\ell_R+912\ell_R
                    \le33\mathcal T\ell_R.
\]
The last step uses \(\mathcal T>912\). The same upper bound
dominates the CM estimate \(854R\). Since
\(\mathcal T\ell_R>128\), the geometric comparison proves
\(h(a)\le67\mathcal T\ell_R\).

As \(M\) is coprime to \(B_*\),
\(\mathcal E(R)=\mathcal E(B_*)\mathcal E(M)\).
The elementary divisor identity gives
\[
 \mathcal E(M)=\sum_{d\mid M}\frac1d
                       \le\sum_{d=1}^M\frac1d\le1+\log M.
\]
Also \(\log(B_*M)\le(\log B_*)(1+\log M)\).
These inequalities prove \eqref{eq:quadratic}.
For \(\eta>0\), the function
\((1+u)^2e^{-\eta u}\), \(u\ge0\), has a finite computable
maximum. For \(\eta=1/3\), its maximum is \(36e^{-5/3}<7\),
which gives the final assertion.
\end{proof}

\section{Scope and verification}

The estimate is uniform over all rational arguments outside the four
finite contacts. It does not restrict the congruence classes of primes
dividing \(M_5(a)\), and it does not assume \(abc\) or Szpiro.
It uses the full normalized compatible representation in the stated
arithmetic source. Merely knowing that \(A\) is one factor of a modular
Jacobian would not justify the multiplicity conclusion.

The finite constant and exponent checks do not prove modularity,
the arithmetic realization, or the isogeny theorem. Those inputs and
the preceding representation and height arguments are separate from
the exact numerical certificate.

\begin{thebibliography}{9}\small
\bibitem{Source}
E.~Glazer and contributors, \emph{Diophantine definition of the integers
in the rationals}, \texttt{manuscripts/quartic\_v9.tex}, commit
\texttt{14e8f601fe0d9173d4a5f084aedc3e07066a7688},
\url{https://github.com/elliotglazer/diophantine-z-q}.
\bibitem{Characters}
J.~Brox and contributors, \emph{Scalar Characters and Contact-Height
Estimates}, \texttt{nebentypus\_and\_contact\_heights\_20261110\_v18.tex},
\url{https://github.com/JoseBrox/diophantine-z-q_improvements}.
\bibitem{Ribet}
K.~A.~Ribet, \emph{Abelian varieties over $\Q$ and modular forms},
Section~3, \url{https://arxiv.org/abs/alg-geom/9208002}.
\bibitem{Faltings}
G.~Faltings, \emph{Endlichkeitss\"atze f\"ur abelsche Variet\"aten
\"uber Zahlk\"orpern}, Inventiones Mathematicae \textbf{73} (1983),
349--366, \url{https://doi.org/10.1007/BF01388432}.
\bibitem{Remond}
G.~R\'emond, \emph{Propri\'et\'es de la hauteur de Faltings},
Annali della Scuola Normale Superiore di Pisa, Classe di Scienze
\textbf{23} (2022), 1589--1596, Corollary~1.2 and the paragraph
following Corollary~1.3,
\url{https://doi.org/10.2422/2036-2145.202010_062}.
\bibitem{GaudronRemond}
\'E.~Gaudron and G.~R\'emond, \emph{Nouveaux th\'eor\`emes d'isog\'enie},
M\'emoires de la Soci\'et\'e Math\'ematique de France
\textbf{176} (2023), Theorem~1.9(1),
\url{https://doi.org/10.24033/msmf.484}.
\end{thebibliography}
\end{document}
